\scanpage{79}
% archive sheet number: 78

\textcircled{3} Étude du cas \(f^*\) pl fidèle. Étude du cas où \(f_*\) commute aux \(\underrightarrow{\lim}\) \unsure{filtrées|finies}. [et aussi, à \unsure{certaines} \(\underrightarrow{\lim}\) ou transforme épim en épim \(\ldots\)] \inserted{-- fam. ép. en fam. ép.}

4) Étude du cas \(f_*\) pl fidèle (Cf \(1^\circ\)))

5) Notion de morphisme \uline{dominant} de topos (\(f_*(\varnothing)=\varnothing\)). \deleted{\unsure{6)|Plus}} Lien avec le fait \(f^*\) \uline{conservatif}. \deleted{\unsure{6)}} (Peut-on dire \unsure{qque} chose sur le cas \(f_*\) conservatif?)

\deleted{6) Donner une caractérisation des foncteurs \unsure{exact|exacts} à gauche \(\unsure{\mathcal U|U}\unsure{-|\longrightarrow}\unsure{\mathcal V|V}\) permettant un recollement à la Cartier-Artin. Cf Deligne)}

% slanted
\margin{\deleted{OK.}}

\deleted{7) Les Morphismes de \uuline{\(\mathcal U\)}-topos \inserted{\(\mathcal U,\mathcal V\)} font une \(\mathcal U\)-cat (i.e. les Hom sont ``petits''). \deleted{et} L'ens des classes, à isom près, de tels morphismes est-il petit ?}

% slanted
\margin{\uuline{non}}

8) a) Théorie des structures : la \unsure{Artin|Notion} des \unsure{'}topos noethériens ?

\diagram{\begin{tikzcd}[ampersand replacement=\&,column sep=huge]
(\mathrm{Top})\arrow[r,"{\text{cat. des fctrs fibres}}","{\scriptscriptstyle >}" marking]\&
(\mathrm{Cat})\arrow[l,bend left=20,"{\widehat C\unsure{\leftarrow|\sim|\mathrel{\leftarrow\!\shortmid}}C}"]
\end{tikzcd}}

\margin{ou y-a-t-il lieu de considérer qu'on s'insère dans une catégorie plus spécifique que (Cat)?}

\inserted{b) Topos où tt ss-objet est \unsure{sommande} (donc où la th. des supports est triviale)}

\inserted{c) Topos engendrés par les ss-objet de l'objet final.}

9) Déterminer la \uline{catégorie} \uline{Hom} de Topos en topos annelés dans certains cas habituels. (On doit obtenir de nombreuses catégories discrètes).

10) Topos modulaires. Catégories fibrées resp. Champs essentiellement représentables \(\ldots\).

\diagram{\begin{tikzcd}[ampersand replacement=\&,column sep=small]
\deleted{\text{11)}}\&\deleted{(\mathrm{Esp.})\unsure{\mathrm{red}}}\arrow[r]\&
\deleted{(\mathrm{Esp.}\ \text{\unsure{sobres}})}\arrow[r]\&
\deleted{(\mathrm{Top})}\arrow[l,bend left=30,"{\deleted{\substack{\text{2-adjoint : espace}\\\text{des foncteurs}\\\text{fibres}}}}"']
\end{tikzcd}}
